% =============================================================================
% APPENDIX HF: PREDICT-HABIT: Habit Formation Mechanism Taxonomy
% =============================================================================
%
% Category: PREDICT
% Module: BCM2_PREDICT_HABIT
% Version: 1.0 (January 2026)
% Status: Draft
% Language: english
%
% This appendix formalizes the taxonomy of habit formation mechanisms,
% distinguishing between learning/experience-good, addiction/reinforcement,
% and inertia/switching-cost models. Provides testable predictions and
% empirical discriminants based on Jones, Molitor & Reif (2025).
%
% Dependencies: AW (CORE-STAGE), AU (CORE-AWARE), BBB (CORE-WHERE)
% Dependents: HHH (METHOD-TOOLKIT), S (PREDICT-MASTER)
%
% =============================================================================

% =============================================================================
% CROSS-REFERENCE STRUCTURE
% =============================================================================
%
% CHAPTER LINKAGE (Required):
% - Primary Chapter: Chapter 13: The Behavioral Change Journey
% - Secondary Chapters: Chapter 14: Health Interventions, Chapter 17: Policy Design
%
% This appendix provides foundation for:
% - Intervention timing optimization (front-load vs. sustained)
% - Mechanism-aware behavioral model design
%
% DEPENDENCIES (what this appendix REQUIRES):
% - Appendix AW (CORE-STAGE): Behavioral Change Journey phases
% - Appendix AU (CORE-AWARE): Awareness function and belief updating
% - Appendix BBB (CORE-WHERE): Parameter values from empirical literature
%
% DEPENDENTS (what REQUIRES this appendix):
% - Appendix HHH (METHOD-TOOLKIT): Intervention design specifications
% - Appendix S (PREDICT-MASTER): Integration with core predictions
%
% =============================================================================

% -----------------------------------------------------------------------------
% HEADER BLOCK (for template compliance)
% -----------------------------------------------------------------------------
\begin{tcolorbox}[colback=gray!5!white, colframe=gray!75!black,
    title=Appendix Metadata]
\textbf{Appendix:} HF (PREDICT-HABIT) \\
\textbf{Category:} PREDICT --- Testable Predictions \\
\textbf{Version:} 1.0 (January 2026) \\
\textbf{Status:} Draft \\
\textbf{Dependencies:} AW (CORE-STAGE), AU (CORE-AWARE), BBB (CORE-WHERE) \\
\textbf{Dependents:} HHH (METHOD-TOOLKIT), S (PREDICT-MASTER)
\end{tcolorbox}

\begin{tcolorbox}[colback=purple!5!white, colframe=purple!75!black,
    title=Chapter Linkage]

\textbf{Primary Chapter:} Chapter 13: The Behavioral Change Journey
\begin{itemize}[nosep]
\item Section 13.3: Phase transitions $\leftrightarrow$ This Appendix Section \ref{sec:hf-theory}
\item Section 13.5: Intervention timing $\leftrightarrow$ This Appendix Section \ref{sec:hf-practice}
\end{itemize}

\textbf{Secondary Chapters:}
\begin{itemize}[nosep]
\item Chapter 14: Health Interventions --- references mechanism selection
\item Chapter 17: Policy Design --- references incentive timing
\end{itemize}

\textbf{Reading Guide:}
\begin{itemize}[nosep]
\item For conceptual overview: Read Chapter 13 first
\item For formal details: This appendix provides complete specification
\item For parameter values: See Appendix BBB (CORE-WHERE)
\end{itemize}

\end{tcolorbox}

\chapter{PREDICT-HABIT: Habit Formation Mechanism Taxonomy}
\label{app:predict-habit}

\begin{abstract}
How do behaviors become habits? This appendix formalizes three competing mechanisms:
(1) learning/experience-good, (2) addiction/reinforcement, and (3) inertia/switching-costs.
Each mechanism generates distinct testable predictions about the temporal pattern of
behavior persistence. We document that Jones, Molitor \& Reif (2025) provides a clean
falsification of addiction models in favor of experience-good learning for annual health
screenings. The key insight: first exposure dominates subsequent exposures in driving
habit formation, with implications for intervention design across domains.
\end{abstract}

\begin{tcolorbox}[colback=blue!5!white, colframe=blue!75!black,
    title=Quick Reference: Habit Formation Mechanisms]

\textbf{Core Question:} What mechanism drives habit formation---and how can we distinguish them empirically?

\textbf{Key Formula:}
\begin{equation}
\Delta P(a_t | a_{t-1}, a_{t-2}) = \begin{cases}
\beta_1 \cdot a_{t-1} + \beta_2 \cdot a_{t-2} & \text{(Addiction: } \beta_1 > \beta_2 \text{)} \\
\beta_1 \cdot a_{t-1} + \beta_2 \cdot a_{t-2} & \text{(Learning: } \beta_1 \approx \beta_2 \text{)} \\
\beta \cdot \mathbb{1}[\exists \tau < t: a_\tau = 1] & \text{(Inertia: any prior action)}
\end{cases}
\end{equation}

\textbf{Key Concepts:}
\begin{itemize}[nosep]
\item \textbf{Experience-Good Learning}: First exposure updates beliefs; subsequent adds little
\item \textbf{Addiction/Reinforcement}: Recent behavior matters more than distant behavior
\item \textbf{Inertia/Switching Costs}: Any participation creates lock-in
\end{itemize}

\textbf{Cross-References:} CORE-STAGE (AW), CORE-AWARE (AU), CORE-WHERE (BBB), PREDICT-MASTER (S), LIT-LEARNING (XIX), METHOD-EVIDENCE (XX)
\end{tcolorbox}

% -----------------------------------------------------------------------------
% APPENDIX SCOPE BOX
% -----------------------------------------------------------------------------

\begin{tcolorbox}[
    colback=orange!5!white,
    colframe=orange!75!black,
    title={\textbf{Appendix Scope: Ziel / In-Scope / Out-of-Scope / Constraints}},
    fonttitle=\bfseries
]

\textbf{Ziel (Objective):}
\begin{itemize}[nosep]
    \item Given that a behavior persists over time, which mechanism is responsible---and how can we tell?
\end{itemize}

\textbf{In-Scope:}
\begin{itemize}[nosep]
    \item The Fundamental Question
    \item Core Theory: The Three Mechanisms
    \item Axioms and Formal Definitions
    \item Formal State-Space Model
    \item Experimental Design and Identification
\end{itemize}

\textbf{Out-of-Scope (delegiert an andere Appendices):}
\begin{itemize}[nosep]
    \item CORE-STAGE $\rightarrow$ Appendix AW
    \item CORE-AWARE $\rightarrow$ Appendix AU
    \item CORE-WHERE $\rightarrow$ Appendix BBB
    \item METHOD-TOOLKIT $\rightarrow$ Appendix HHH
    \item PREDICT-MASTER $\rightarrow$ Appendix S
\end{itemize}

\textbf{Constraints (Anwendungsgrenzen):}
\begin{itemize}[nosep]
    \item [TODO: Define when this appendix does NOT apply]
    \item [TODO: Define required prerequisites]
    \item [TODO: Define known limitations]
\end{itemize}

\textbf{Lieferobjekte (Deliverables):}
\begin{itemize}[nosep]
    \item 9 Table(s)
    \item 22 Equation(s)
    \item 5 Axiom(s)
    \item Worked Example(s)
\end{itemize}

\end{tcolorbox}




%%%%%%%%%%%%%%%%%%%%%%%%%%%%%%%%%%%%%%%%%%%%%%%%%%%%%%%%%%%%%%%%%%%%%%%%%%%%%%%
\section{The Fundamental Question}
\label{sec:hf-fundamental}
%%%%%%%%%%%%%%%%%%%%%%%%%%%%%%%%%%%%%%%%%%%%%%%%%%%%%%%%%%%%%%%%%%%%%%%%%%%%%%%

\subsection{Why This Matters}

Habit formation is central to behavioral economics and intervention design. If we understand
\textit{why} behaviors persist, we can design more effective interventions. The key question
is not just whether habits form, but \textit{through which mechanism}---because different
mechanisms imply radically different optimal intervention strategies.

Consider three intervention philosophies:
\begin{itemize}
\item \textbf{Front-load}: Invest heavily in first-time experience quality and uptake
\item \textbf{Sustain}: Provide ongoing reinforcement through repeated incentives
\item \textbf{Remove friction}: Eliminate switching costs and barriers to re-engagement
\end{itemize}

The correct choice depends entirely on which mechanism drives habit formation in the target domain.

\subsection{The Core Problem}

Most habit research cannot distinguish between mechanisms because it lacks:
\begin{enumerate}
\item Multiple periods of observation
\item Exogenous variation in behavior at each period
\item Ability to separate effects of initial vs. recent past behavior
\end{enumerate}

\begin{tcolorbox}[colback=red!5!white, colframe=red!75!black,
    title=The Fundamental Question]
\textbf{Given that a behavior persists over time, which mechanism is responsible---and how can we tell?}
\end{tcolorbox}


%%%%%%%%%%%%%%%%%%%%%%%%%%%%%%%%%%%%%%%%%%%%%%%%%%%%%%%%%%%%%%%%%%%%%%%%%%%%%%%
\section{Core Theory: The Three Mechanisms}
\label{sec:hf-theory}
%%%%%%%%%%%%%%%%%%%%%%%%%%%%%%%%%%%%%%%%%%%%%%%%%%%%%%%%%%%%%%%%%%%%%%%%%%%%%%%

\subsection{Mechanism 1: Experience-Good Learning}

An \textbf{experience good} is a product or service whose value can only be assessed through
consumption. Prior to first experience, agents have uncertain beliefs about costs and benefits.

\begin{tcolorbox}[colback=gray!5!white, colframe=gray!75!black,
    title=Definition: Experience-Good Learning]
Let $\mu_0$ be prior belief about net benefit of action $a$, and $\sigma^2_0$ the prior variance.
After first experience, beliefs update via:
\begin{equation}
\mu_1 = \frac{\sigma^2_0}{\sigma^2_0 + \sigma^2_\epsilon} \cdot x + \frac{\sigma^2_\epsilon}{\sigma^2_0 + \sigma^2_\epsilon} \cdot \mu_0
\end{equation}
where $x$ is realized benefit and $\sigma^2_\epsilon$ is experience noise.

\textbf{Key property:} Second experience provides minimal additional information if $\sigma^2_1 \ll \sigma^2_0$.
\end{tcolorbox}

\textbf{Prediction:} First exposure has large effect on subsequent behavior; additional exposures add little.

\textbf{EBF Connection:} Maps to CORE-AWARE (AU) via belief updating function $A(\cdot)$.
The awareness level $A_i$ jumps discontinuously at first exposure.

\subsection{Mechanism 2: Addiction/Reinforcement}

In addiction models, past consumption directly increases the marginal utility of current
consumption. The ``stock'' of past behavior creates physiological or psychological dependence.

\begin{tcolorbox}[colback=gray!5!white, colframe=gray!75!black,
    title=Definition: Addiction/Reinforcement]
Let $S_t$ be the ``stock'' of past behavior, with accumulation:
\begin{equation}
S_t = \delta S_{t-1} + a_{t-1}, \quad \delta \in (0,1)
\end{equation}
Utility from action depends on stock:
\begin{equation}
U(a_t, S_t) = u(a_t) + \gamma \cdot S_t \cdot a_t
\end{equation}

\textbf{Key property:} More recent behavior has higher weight due to depreciation $\delta < 1$.
\end{tcolorbox}

\textbf{Prediction:} Recent past behavior matters more than distant past behavior.

\textbf{EBF Connection:} Maps to CORE-HOW (B) via complementarity $\gamma$ between current
action and accumulated stock.

\subsection{Mechanism 3: Inertia/Switching Costs}

Switching costs create path dependence: any prior participation reduces the cost of
future participation, regardless of recency.

\begin{tcolorbox}[colback=gray!5!white, colframe=gray!75!black,
    title=Definition: Inertia/Switching Costs]
Let $c(a_t | H_t)$ be the cost of action given history $H_t$:
\begin{equation}
c(a_t = 1 | H_t) = \begin{cases}
c_0 & \text{if } a_\tau = 0 \text{ for all } \tau < t \\
c_0 - \kappa & \text{if } \exists \tau < t: a_\tau = 1
\end{cases}
\end{equation}
where $\kappa > 0$ is the switching cost reduction from any prior experience.

\textbf{Key property:} Timing of prior experience is irrelevant; only existence matters.
\end{tcolorbox}

\textbf{Prediction:} Any prior participation increases future participation equally, regardless of when.

\textbf{EBF Connection:} Maps to CORE-WHEN (V) via temporal context $\Psi_T$, but
predicts context-independence within the experienced state.


%%%%%%%%%%%%%%%%%%%%%%%%%%%%%%%%%%%%%%%%%%%%%%%%%%%%%%%%%%%%%%%%%%%%%%%%%%%%%%%
\section{Axioms and Formal Definitions}
\label{sec:hf-axioms}
%%%%%%%%%%%%%%%%%%%%%%%%%%%%%%%%%%%%%%%%%%%%%%%%%%%%%%%%%%%%%%%%%%%%%%%%%%%%%%%

\begin{tcolorbox}[colback=gray!5!white, colframe=gray!75!black,
    title=Axiom HF-1: Mechanism Exclusivity (Local)]
\textbf{Statement:} For a given behavior-domain pair $(a, D)$, exactly one primary mechanism
governs habit formation:
\begin{equation}
\mathcal{M}(a, D) \in \{\text{LEARN}, \text{ADDICT}, \text{INERTIA}\}
\end{equation}

\textbf{Interpretation:} While multiple mechanisms may contribute, one dominates empirically.

\textbf{Implication:} Optimal intervention strategy is determined by $\mathcal{M}(a, D)$.
\end{tcolorbox}

\begin{tcolorbox}[colback=gray!5!white, colframe=gray!75!black,
    title=Axiom HF-2: Temporal Discriminant]
\textbf{Statement:} The three mechanisms generate distinct signatures in the coefficients
$(\beta_1, \beta_2)$ of the regression:
\begin{equation}
P(a_t = 1) = \alpha + \beta_1 \cdot a_{t-1} + \beta_2 \cdot a_{t-2} + \epsilon_t
\end{equation}

\textbf{Discriminant Rules:}
\begin{align}
\text{ADDICT:} & \quad \beta_1 > \beta_2 > 0 \\
\text{LEARN:} & \quad \beta_1 \approx \beta_2 > 0 \text{ and } \beta_2 \text{ large from first exposure} \\
\text{INERTIA:} & \quad \beta_1 + \beta_2 \approx \text{constant}
\end{align}

\textbf{Implication:} Multi-period experimental data can falsify mechanisms.
\end{tcolorbox}

\begin{tcolorbox}[colback=gray!5!white, colframe=gray!75!black,
    title=Axiom HF-3: Front-Loading Principle]
\textbf{Statement:} If $\mathcal{M}(a, D) = \text{LEARN}$, then:
\begin{equation}
\frac{\partial V}{\partial I_1} > \frac{\partial V}{\partial I_t} \quad \forall t > 1
\end{equation}
where $V$ is total long-run behavior and $I_t$ is incentive at time $t$.

\textbf{Interpretation:} Investment in first-time uptake has higher ROI than later incentives.

\textbf{Implication:} Ignoring habit formation undervalues initial interventions by $\approx 70\%$
(Jones et al. 2025 estimate).
\end{tcolorbox}

\begin{tcolorbox}[colback=gray!5!white, colframe=gray!75!black,
    title=Axiom HF-4: Belief Asymmetry]
\textbf{Statement:} For experience goods, expected uncertainty reduction is:
\begin{equation}
\mathbb{E}[\Delta \sigma^2 | \text{1st exposure}] \gg \mathbb{E}[\Delta \sigma^2 | \text{2nd exposure}]
\end{equation}

\textbf{Interpretation:} First exposure resolves most belief uncertainty.

\textbf{Implication:} Mechanism LEARN predicts asymmetric effects between first and subsequent exposures.
\end{tcolorbox}


%%%%%%%%%%%%%%%%%%%%%%%%%%%%%%%%%%%%%%%%%%%%%%%%%%%%%%%%%%%%%%%%%%%%%%%%%%%%%%%
\section{Formal State-Space Model}
\label{sec:hf-statespace}
%%%%%%%%%%%%%%%%%%%%%%%%%%%%%%%%%%%%%%%%%%%%%%%%%%%%%%%%%%%%%%%%%%%%%%%%%%%%%%%

The following minimal model captures the essential dynamics of experience-good learning
with an \textbf{absorbing state}---a parsimonious formalization that nests the key empirical
predictions.

\subsection{The Absorbing-State Model}

\begin{tcolorbox}[colback=orange!5!white, colframe=orange!75!black,
    title=Definition: Minimal State-Space Model for Experience-Good Learning]

\textbf{State Variable:}
\begin{equation}
H_{it} \in \{0, 1\} \quad \text{(binary: has individual } i \text{ ever participated?)}
\end{equation}

\textbf{State Transition (Absorbing):}
\begin{equation}
H_{i,t+1} = \max(H_{it}, S_{it})
\end{equation}
where $S_{it} \in \{0,1\}$ is the participation decision at time $t$.

\textbf{Utility:}
\begin{equation}
U_{it} = u(H_{it}) - c + \varepsilon_{it}, \quad \text{with } u(1) > u(0)
\end{equation}
where $c$ is the participation cost and $\varepsilon_{it}$ is an i.i.d. preference shock.

\textbf{Decision Rule:}
\begin{equation}
S_{it} = \mathbf{1}\{U_{it} + Z_{it} \geq 0\}
\end{equation}
where $Z_{it}$ is the (randomly assigned) financial incentive.

\end{tcolorbox}

\subsection{Key Model Properties}

\begin{enumerate}
\item \textbf{Non-additive dynamics:} The state $H$ does not accumulate---it switches once.

\item \textbf{Absorbing state:} Once $H_{it} = 1$, the state remains $H_{i,t'} = 1$ for all $t' > t$.

\item \textbf{No depreciation:} Unlike addiction models, there is no decay parameter $\delta$.

\item \textbf{Stable long-run persistence:} The utility gain $u(1) - u(0)$ is permanent.
\end{enumerate}

\subsection{Contrast with Addiction Model}

The standard myopic addiction model (Becker \& Murphy 1988) implies:
\begin{equation}
a_t = k + \alpha S_t + \pi p_t, \quad \text{with } S_{t+1} = (1-d) S_t + a_t
\end{equation}

For a three-period setting, this yields:
\begin{equation}
a_3 = \alpha \cdot a_2 + \alpha(1-d) \cdot a_1 + \ldots
\end{equation}

\textbf{Central prediction:} Effect of $a_2$ on $a_3$ \textbf{exceeds} effect of $a_1$ on $a_3$
(recent consumption dominates).

$\Rightarrow$ \textbf{Contradicted by data:} Jones et al. find $\theta_1 \gg \theta_2 \approx 0$.

\subsection{Interpretation as Expectation Update}

The absorbing-state model has an intuitive cognitive interpretation:

\begin{tcolorbox}[colback=blue!5!white, colframe=blue!75!black,
    title=Cognitive Interpretation: One-Shot Belief Revision]

\textbf{Prior beliefs (before first experience):}
\begin{equation}
\mathbb{E}_0[U] < U_{\text{true}} \quad \text{(systematic underestimation of net benefit)}
\end{equation}

\textbf{First experience triggers belief update:}
\begin{equation}
\mathbb{E}_1[U] = \mathbb{E}[U \mid \text{experience}] \approx U_{\text{true}}
\end{equation}

\textbf{Subsequent experiences add no information:}
\begin{equation}
\mathbb{E}_t[U] \approx \mathbb{E}_1[U] \quad \forall t \geq 2
\end{equation}

\textbf{Result:} Behavior is persistent but \textbf{not self-reinforcing}.
The system reaches a new steady state after one transition.

\end{tcolorbox}

This interpretation explains why:
\begin{itemize}
\item First participation has a large, permanent effect ($\theta_1 \approx 0.35$)
\item Second participation adds almost nothing ($\theta_2 \approx 0.03$, n.s.)
\item Effects do not decay over time (no depreciation in beliefs)
\end{itemize}

\subsection{From Model to Intuition: ``The First Impression''}

\begin{tcolorbox}[colback=yellow!5!white, colframe=yellow!75!black,
    title=Intuition Box: Why ``First Impressions Matter'' Is Right---But Not Why You Think]

\textbf{The everyday intuition:}
\begin{quote}
\textit{``Der erste Eindruck ist der wichtigste.''} \\
(``The first impression is the most important.'')
\end{quote}

\textbf{This is empirically correct}, but the mechanism is \textbf{not}:
\begin{itemize}[nosep]
\item[$\times$] \textbf{Primacy bias} (cognitive distortion) --- the paper measures rational learning
\item[$\times$] \textbf{Emotional overwhelm} (``wow effect'') --- it's information updating, not affect
\item[$\times$] \textbf{Diminishing marginal utility} --- there is no gradual decline, but a discrete jump
\end{itemize}

\textbf{The correct interpretation:}
\begin{quote}
\textit{``The first impression matters most---because it permanently changes the expectation state.''}
\end{quote}

\textbf{Formally:} Not $U_t = U_{t-1} + \delta$ (accumulation), but:
\begin{equation}
U_t = \begin{cases}
u_0 & \text{if } H_t = 0 \text{ (never experienced)} \\
u_1 & \text{if } H_t = 1 \text{ (experienced once)}
\end{cases}
\quad \text{with } u_1 > u_0
\end{equation}

$\Rightarrow$ A single jump, not an accumulation process.

\end{tcolorbox}

\begin{tcolorbox}[colback=green!5!white, colframe=green!75!black,
    title=Three Formulations for Different Audiences]

\textbf{(A) Economic-strict} (for models, papers, policy memos):
\begin{quote}
\textit{``First contact determines the long-run behavioral state; subsequent contacts barely change it.''}
\end{quote}

\textbf{(B) Bridge formulation} (for management, ESG, strategy):
\begin{quote}
\textit{``The first impression is decisive because it permanently sets expectations.''}
\end{quote}

\textbf{(C) Pointed} (for executive summaries, talks):
\begin{quote}
\textit{``Habit is not created by repetition, but by the first experience.''}
\end{quote}

\end{tcolorbox}

\subsection{Domain Applications of the First-Impression Principle}

\begin{table}[htbp]
\centering
\caption{Applying the First-Impression Principle Across Domains}
\label{tab:hf-first-impression}
\renewcommand{\arraystretch}{1.4}
\begin{tabular}{@{}p{2.5cm}p{4cm}p{5.5cm}@{}}
\toprule
\textbf{Domain} & \textbf{Application} & \textbf{Implication} \\
\midrule
\textbf{Products / UX} & Onboarding, trial versions, first use & Design first-use experience, not retention loops \\
\textbf{ESG / Sustainability} & CO$_2$ tracking, energy feedback, sustainable defaults & Subsidize entry, not repetition \\
\textbf{Public Health} & Vaccinations, screenings, preventive care & One well-designed initial push dominates many small follow-ups \\
\textbf{HR / Workplace} & Employee wellness programs, training & Invest in first participation quality \\
\textbf{Finance} & Retirement enrollment, savings programs & Default into first contribution, not ongoing nudges \\
\bottomrule
\end{tabular}
\end{table}

\textbf{Key insight:} The first impression matters not because it is \textit{emotional}, but because it is \textit{informative}.


%%%%%%%%%%%%%%%%%%%%%%%%%%%%%%%%%%%%%%%%%%%%%%%%%%%%%%%%%%%%%%%%%%%%%%%%%%%%%%%
\section{Experimental Design and Identification}
\label{sec:hf-identification}
%%%%%%%%%%%%%%%%%%%%%%%%%%%%%%%%%%%%%%%%%%%%%%%%%%%%%%%%%%%%%%%%%%%%%%%%%%%%%%%

\subsection{The Jones-Molitor-Reif Design}

The experimental setting provides an ideal \textbf{identification laboratory}:

\begin{table}[htbp]
\centering
\caption{Experimental Design Features}
\label{tab:hf-design}
\renewcommand{\arraystretch}{1.3}
\begin{tabular}{@{}ll@{}}
\toprule
\textbf{Feature} & \textbf{Specification} \\
\midrule
Sample & 4,799 employees (two cohorts) \\
Duration & Up to 3 screening years \\
Randomization & Individual-level, annual \\
Key innovation & \textbf{Rerandomization} each year \\
Incentive independence & $Z_{it} \perp Z_{i,t-1}, S_{i,t-1}$ \\
\bottomrule
\end{tabular}
\end{table}

\subsection{Two Levels of Analysis}

The authors distinguish between:

\textbf{(A) Reduced-Form Persistence:} Do past \textit{incentives} affect current behavior?
\begin{equation}
S_{it} = \alpha + \sum_{\tau < t} \beta_\tau Z_{i\tau} + \beta_t Z_{it} + \varepsilon_{it}
\end{equation}

\textbf{(B) Structural Habit Formation:} Does past \textit{behavior} causally affect current behavior?
\begin{equation}
S_{it} = \alpha + \sum_{\tau < t} \theta_\tau S_{i\tau} + \gamma Z_{it} + u_{it}
\end{equation}

The structural parameters $\theta_\tau$ are the habit formation effects of interest.

\subsection{Instrumental Variables Strategy}

\textbf{Problem:} $S_{i\tau}$ is endogenous (correlated with stable preferences, unobserved heterogeneity).

\textbf{Solution:} Use randomly assigned incentives $Z_{i\tau}$ as instruments.

\begin{tcolorbox}[colback=gray!5!white, colframe=gray!75!black,
    title=IV Identification]
Under random assignment, the causal effects are identified as:
\begin{align}
\theta_1^{\text{causal}} &= \frac{\text{Cov}(S_3, Z_1)}{\text{Cov}(S_1, Z_1)} \cdot \frac{1}{P(S_2 = 1 | S_1 = 1)} \\[0.5em]
\theta_2^{\text{causal}} &= \frac{\text{Cov}(S_3, Z_2 | S_1)}{\text{Cov}(S_2, Z_2 | S_1)}
\end{align}

\textbf{Key:} Rerandomization at $t=2$ provides variation in $S_2$ that is independent of $S_1$
conditional on the randomization.
\end{tcolorbox}


%%%%%%%%%%%%%%%%%%%%%%%%%%%%%%%%%%%%%%%%%%%%%%%%%%%%%%%%%%%%%%%%%%%%%%%%%%%%%%%
\section{Empirical Results}
\label{sec:hf-results}
%%%%%%%%%%%%%%%%%%%%%%%%%%%%%%%%%%%%%%%%%%%%%%%%%%%%%%%%%%%%%%%%%%%%%%%%%%%%%%%

\begin{proposition}[Jones-Molitor-Reif Identification]
In a three-period randomized experiment with rerandomized incentives, the causal effects
of first ($\beta^{causal}_1$) and second ($\beta^{causal}_2$) exposure can be separately identified
using instrumental variables.
\end{proposition}

\begin{proof}
Let $Z_t$ be the incentive assignment at time $t$. Under random assignment:
\begin{align}
\beta^{causal}_1 &= \frac{\text{Cov}(a_3, Z_1)}{\text{Cov}(a_1, Z_1)} \cdot \frac{1}{P(a_2 = 1 | a_1 = 1)} \\
\beta^{causal}_2 &= \frac{\text{Cov}(a_3, Z_2 | a_1)}{\text{Cov}(a_2, Z_2 | a_1)}
\end{align}
Rerandomization at $t=2$ provides independent variation conditional on $a_1$.
\end{proof}

\subsection{Jones et al. (2025) Findings}

\begin{table}[htbp]
\centering
\caption{Habit Formation in Annual Health Screenings (Jones, Molitor \& Reif 2025)}
\label{tab:hf-jmr-results}
\renewcommand{\arraystretch}{1.4}
\begin{tabular}{@{}lccl@{}}
\toprule
\textbf{Effect} & \textbf{Estimate} & \textbf{95\% CI} & \textbf{Interpretation} \\
\midrule
First exposure $\rightarrow$ Year 2 & +32.4 pp & [28.1, 36.7] & Strong habit formation \\
First exposure $\rightarrow$ Year 3 & +36.0 pp & [30.8, 41.2] & No decay over time \\
Second exposure $\rightarrow$ Year 3 & +3.2 pp & [-2.4, 8.8] & Minimal additional effect \\
\midrule
Relative increase (Year 2) & 84\% & --- & vs. baseline 38.6\% \\
Relative increase (Year 3) & 90\% & --- & vs. baseline 40.0\% \\
\bottomrule
\end{tabular}
\end{table}

\begin{tcolorbox}[colback=green!5!white, colframe=green!75!black,
    title=Key Finding: Learning Mechanism Confirmed]

\textbf{Pattern observed:}
\begin{itemize}[nosep]
\item First exposure: Large, persistent effect (32-36 pp)
\item Second exposure: Small, statistically insignificant effect (3 pp)
\item No decay: Year 3 effect $\geq$ Year 2 effect
\end{itemize}

\textbf{Mechanism verdict:}
\begin{itemize}[nosep]
\item[$\times$] \textbf{Addiction/Reinforcement}: REJECTED (recent $\not>$ distant)
\item[$\checkmark$] \textbf{Experience-Good Learning}: SUPPORTED (first exposure dominates)
\item[?] \textbf{Inertia}: Partially consistent, but does not explain asymmetry
\end{itemize}

\end{tcolorbox}


%%%%%%%%%%%%%%%%%%%%%%%%%%%%%%%%%%%%%%%%%%%%%%%%%%%%%%%%%%%%%%%%%%%%%%%%%%%%%%%
\section{Integration with Other Appendices}
\label{sec:hf-integration}
%%%%%%%%%%%%%%%%%%%%%%%%%%%%%%%%%%%%%%%%%%%%%%%%%%%%%%%%%%%%%%%%%%%%%%%%%%%%%%%

\subsection{Connection to CORE-STAGE (AW)}

The Behavioral Change Journey (BCJ) framework identifies phases:
Precontemplation $\rightarrow$ Contemplation $\rightarrow$ Preparation $\rightarrow$ Action $\rightarrow$ Maintenance.

\begin{tcolorbox}[colback=yellow!5!white, colframe=yellow!75!black,
    title=Integration: PREDICT-HABIT $\leftrightarrow$ CORE-STAGE]
\textbf{What this appendix provides:}
\begin{itemize}[nosep]
\item Mechanism for Action $\rightarrow$ Maintenance transition
\item Asymmetric importance of first Action vs. subsequent Actions
\end{itemize}

\textbf{What it requires:}
\begin{itemize}[nosep]
\item BCJ phase classification $\varphi$
\item Phase transition probabilities $P(\varphi_{t+1} | \varphi_t, a_t)$
\end{itemize}

\textbf{Key Insight:} For learning-amenable behaviors, the critical transition is
Pre-Action $\rightarrow$ First Action, not subsequent maintenance.
\end{tcolorbox}

\subsection{Connection to CORE-AWARE (AU)}

The Awareness function $A(\cdot)$ captures subjective perception of utility.

\begin{tcolorbox}[colback=yellow!5!white, colframe=yellow!75!black,
    title=Integration: PREDICT-HABIT $\leftrightarrow$ CORE-AWARE]
\textbf{What this appendix provides:}
\begin{itemize}[nosep]
\item Discontinuous jump in $A_i$ at first exposure
\item Quantification: $\Delta A^{(1)} \gg \Delta A^{(k)}$ for $k > 1$
\end{itemize}

\textbf{What it requires:}
\begin{itemize}[nosep]
\item Prior belief structure $(\mu_0, \sigma^2_0)$
\item Experience signal precision $\sigma^2_\epsilon$
\end{itemize}

\textbf{Mathematical Interface:}
\begin{equation}
A^{(k)}_i = A^{(0)}_i + \sum_{j=1}^{k} \Delta A^{(j)}_i, \quad \text{where } \Delta A^{(1)} \approx 0.3\text{--}0.4
\end{equation}
\end{tcolorbox}

\subsection{Connection to CORE-WHERE (BBB)}

\begin{tcolorbox}[colback=yellow!5!white, colframe=yellow!75!black,
    title=Integration: PREDICT-HABIT $\leftrightarrow$ CORE-WHERE]
\textbf{Parameter contributions from this appendix:}

\begin{table}[H]
\centering
\begin{tabular}{@{}llll@{}}
\toprule
\textbf{Parameter} & \textbf{Value} & \textbf{Source} & \textbf{Domain} \\
\midrule
$\beta^{\text{LEARN}}_1$ (first exposure) & 0.32--0.36 & Jones et al. 2025 & Health screening \\
$\beta^{\text{LEARN}}_2$ (second exposure) & 0.03 (n.s.) & Jones et al. 2025 & Health screening \\
Decay rate (2 years) & $\approx 0$ & Jones et al. 2025 & Health screening \\
Undervaluation factor & 70\% & Jones et al. 2025 & Cost-effectiveness \\
\bottomrule
\end{tabular}
\end{table}
\end{tcolorbox}


%%%%%%%%%%%%%%%%%%%%%%%%%%%%%%%%%%%%%%%%%%%%%%%%%%%%%%%%%%%%%%%%%%%%%%%%%%%%%%%
\section{Worked Example: Corporate Wellness Program}
\label{sec:hf-example}
%%%%%%%%%%%%%%%%%%%%%%%%%%%%%%%%%%%%%%%%%%%%%%%%%%%%%%%%%%%%%%%%%%%%%%%%%%%%%%%

\subsection{Scenario: Annual Health Screening Incentive Design}

\paragraph{Setup.} A company wants to maximize long-term health screening participation.
Budget: \$100 per employee over 3 years. Current baseline participation: 40\%.

\paragraph{Strategy Comparison.}

\begin{table}[htbp]
\centering
\caption{Three Incentive Strategies}
\label{tab:hf-strategies}
\renewcommand{\arraystretch}{1.3}
\begin{tabular}{@{}lcccl@{}}
\toprule
\textbf{Strategy} & \textbf{Year 1} & \textbf{Year 2} & \textbf{Year 3} & \textbf{Mechanism Assumed} \\
\midrule
Uniform & \$33 & \$33 & \$33 & None (naive) \\
Front-loaded & \$75 & \$25 & \$0 & Learning \\
Back-loaded & \$0 & \$25 & \$75 & Addiction \\
\bottomrule
\end{tabular}
\end{table}

\paragraph{Application.} Using Jones et al. parameters ($\beta_1 = 0.35$, incentive elasticity = 0.15 per \$25):

\textbf{Front-loaded Strategy:}
\begin{align}
P_1 &= 0.40 + 0.15 \times 3 = 0.85 \\
P_2 &= 0.40 + 0.35 \times 0.85 + 0.15 \times 1 = 0.85 \\
P_3 &= 0.40 + 0.35 \times 0.85 + 0.03 \times 0.85 = 0.72
\end{align}

\textbf{Uniform Strategy:}
\begin{align}
P_1 &= 0.40 + 0.15 \times 1.32 = 0.60 \\
P_2 &= 0.40 + 0.35 \times 0.60 + 0.15 \times 1.32 = 0.81 \\
P_3 &= 0.40 + 0.35 \times 0.60 + 0.03 \times 0.81 + 0.15 \times 1.32 = 0.83
\end{align}

\paragraph{Result.} Total screenings over 3 years:
\begin{itemize}
\item Front-loaded: $0.85 + 0.85 + 0.72 = 2.42$ per employee
\item Uniform: $0.60 + 0.81 + 0.83 = 2.24$ per employee
\end{itemize}

\textbf{Front-loading outperforms by 8\% with same budget.}

\paragraph{Discussion.} The key insight is that the learning mechanism creates a
multiplier effect: high Year 1 participation generates ``free'' participation in
subsequent years through habit formation. The naive uniform strategy underinvests
in the most valuable period.


%%%%%%%%%%%%%%%%%%%%%%%%%%%%%%%%%%%%%%%%%%%%%%%%%%%%%%%%%%%%%%%%%%%%%%%%%%%%%%%
\section{Implications for Practice}
\label{sec:hf-practice}
%%%%%%%%%%%%%%%%%%%%%%%%%%%%%%%%%%%%%%%%%%%%%%%%%%%%%%%%%%%%%%%%%%%%%%%%%%%%%%%

\begin{tcolorbox}[colback=green!5!white, colframe=green!75!black,
    title=Practical Implications]
\begin{enumerate}
\item \textbf{Front-load incentives for learning-amenable behaviors:}
   Initial uptake investment has 70\%+ higher long-run ROI than sustained incentives.

\item \textbf{Diagnose the mechanism before designing interventions:}
   Use multi-period data to test whether recent or distant past behavior matters more.

\item \textbf{Optimize first-time experience quality:}
   If learning mechanism dominates, experience quality affects permanent beliefs.

\item \textbf{Do not expect compounding returns from repeated incentives:}
   Unlike addiction, repeated ``doses'' do not generate increasing persistence.

\item \textbf{Apply beyond health:}
   Learning mechanism likely relevant for: vaccinations, preventive care, onboarding,
   ESG behaviors, retirement planning---any low-frequency behavior with uncertain value.
\end{enumerate}
\end{tcolorbox}

\subsection{Domain-Mechanism Mapping}

\begin{table}[htbp]
\centering
\caption{Mechanism Discrimination: Formal Predictions vs. Empirical Evidence}
\label{tab:hf-mechanism-discrimination}
\renewcommand{\arraystretch}{1.4}
\begin{tabular}{@{}p{3cm}p{4.5cm}p{3cm}c@{}}
\toprule
\textbf{Mechanism} & \textbf{Formal Prediction} & \textbf{Jones et al. Data} & \textbf{Verdict} \\
\midrule
Addiction / Reinforcement & Recent $>$ distant: $\theta_1 > \theta_2 > 0$ & $\theta_1 \approx 0.35$, $\theta_2 \approx 0.03$ & $\times$ Rejected \\
Experience-Good Learning & First dominates: $\theta_1 \gg \theta_2 \approx 0$ & Exactly observed & $\checkmark$ Supported \\
Inertia / Switching Costs & Any prior sufficient: depends on $\exists$ not timing & Partial match & ? Inconclusive \\
\bottomrule
\end{tabular}
\end{table}

\begin{table}[htbp]
\centering
\caption{Hypothesized Dominant Mechanisms by Domain}
\label{tab:hf-domain-mechanisms}
\renewcommand{\arraystretch}{1.3}
\begin{tabular}{@{}llll@{}}
\toprule
\textbf{Domain} & \textbf{Behavior} & \textbf{Likely Mechanism} & \textbf{Evidence} \\
\midrule
Health & Annual screening & Learning & Jones et al. 2025 \\
Health & Handwashing & Reinforcement & Hussam et al. 2022 \\
Finance & Retirement enrollment & Inertia & Madrian \& Shea 2001 \\
Exercise & Gym attendance & Mixed & Charness \& Gneezy 2009 \\
Digital & App engagement & Addiction & Behavioral design lit. \\
\bottomrule
\end{tabular}
\end{table}


%%%%%%%%%%%%%%%%%%%%%%%%%%%%%%%%%%%%%%%%%%%%%%%%%%%%%%%%%%%%%%%%%%%%%%%%%%%%%%%
\section{Summary}
\label{sec:hf-summary}
%%%%%%%%%%%%%%%%%%%%%%%%%%%%%%%%%%%%%%%%%%%%%%%%%%%%%%%%%%%%%%%%%%%%%%%%%%%%%%%

\begin{tcolorbox}[colback=green!10!white, colframe=green!75!black,
    title=\textbf{Appendix HF: Summary}]

\textbf{Core Question:} What mechanism drives habit formation---and how can we distinguish them empirically?

\textbf{Main Contribution:}
\begin{itemize}[nosep]
    \item Taxonomy of three habit formation mechanisms with testable predictions
    \item Formal \textbf{absorbing-state model} for experience-good learning:
    $H_{i,t+1} = \max(H_{it}, S_{it})$
    \item Clean falsification: Addiction predicts $\theta_1 > \theta_2$; data show $\theta_1 \gg \theta_2 \approx 0$
    \item Integration of Jones et al. (2025) with IV identification strategy
\end{itemize}

\textbf{Key Outputs:}
\begin{itemize}[nosep]
    \item Axioms HF-1 through HF-4 for mechanism identification
    \item Structural parameters: $\theta_1 = 0.32$--$0.36$, $\theta_2 \approx 0.03$ (n.s.)
    \item Cognitive interpretation: One-shot belief revision $\mathbb{E}_0[U] \to \mathbb{E}_1[U]$
    \item Front-loading principle for learning-amenable behaviors (70\% ROI gain)
\end{itemize}

\textbf{Key Insight for Formal Models:}
\begin{quote}
\textit{Habit formation $\neq$ reinforcement. For experience goods, it is a discrete
learning step with an absorbing state---not a gradual accumulation process.
The state space is smaller than addiction models suggest: binary, not continuous.}
\end{quote}

\textbf{Integration:} Connects to CORE-STAGE (BCJ: discrete phase transition),
CORE-AWARE (belief updating with $\Delta A^{(1)} \gg \Delta A^{(k)}$),
and CORE-WHERE (empirical parameters from IV estimation).

\end{tcolorbox}


%%%%%%%%%%%%%%%%%%%%%%%%%%%%%%%%%%%%%%%%%%%%%%%%%%%%%%%%%%%%%%%%%%%%%%%%%%%%%%%
\section{Glossary of Symbols}
\label{sec:hf-glossary}
%%%%%%%%%%%%%%%%%%%%%%%%%%%%%%%%%%%%%%%%%%%%%%%%%%%%%%%%%%%%%%%%%%%%%%%%%%%%%%%

\begin{tcolorbox}[colback=blue!5!white, colframe=blue!75!black]
\textbf{Central Reference:} For complete symbol definitions, see
\textbf{Appendix G} (REF-GLOSSARY).

This section lists only symbols introduced or specialized in this appendix.
\end{tcolorbox}

\vspace{0.5em}

\begin{table}[htbp]
\centering
\caption{Symbols Introduced in Appendix HF}
\label{tab:hf-symbols}
\renewcommand{\arraystretch}{1.3}
\begin{tabular}{@{}clp{5.5cm}c@{}}
\toprule
\textbf{Symbol} & \textbf{Name} & \textbf{Definition} & \textbf{Also in G?} \\
\midrule
$\mathcal{M}(a,D)$ & Mechanism & Dominant habit formation mechanism for action $a$ in domain $D$ & NEW \\
$H_{it}$ & Experience state & Binary: has individual $i$ ever participated? (absorbing) & NEW \\
$\theta_1, \theta_2$ & Structural habit effects & Causal effect of past behavior on current (IV-estimated) & NEW \\
$\beta_1, \beta_2$ & Reduced-form coefficients & Effect of past incentives on current behavior & NEW \\
$S_t$ & Addiction stock & Accumulated past behavior with depreciation & NEW \\
$Z_{it}$ & Incentive assignment & Randomly assigned financial incentive at time $t$ & NEW \\
$\kappa$ & Switching cost reduction & Cost reduction from prior experience & NEW \\
$\mu_0, \sigma^2_0$ & Prior beliefs & Mean and variance of value beliefs pre-experience & \checkmark \\
$\mathbb{E}_t[U]$ & Expected utility & Beliefs about utility at time $t$ (after $t$ experiences) & NEW \\
\bottomrule
\end{tabular}
\end{table}


%%%%%%%%%%%%%%%%%%%%%%%%%%%%%%%%%%%%%%%%%%%%%%%%%%%%%%%%%%%%%%%%%%%%%%%%%%%%%%%
\section{Critical Foundations and Objections}
\label{sec:hf-foundations}
%%%%%%%%%%%%%%%%%%%%%%%%%%%%%%%%%%%%%%%%%%%%%%%%%%%%%%%%%%%%%%%%%%%%%%%%%%%%%%%

\subsection{Objection 1: External Validity}

\textbf{Objection:} Jones et al. studies only health screenings at one employer.
Results may not generalize to other behaviors or contexts.

\textbf{Response:} The theoretical prediction (learning mechanism implies first exposure dominance)
is general. Health screening is a particularly clean test case because it is:
(a) annual and low-frequency, (b) has uncertain value to participants,
(c) can be cleanly randomized. The mechanism taxonomy applies broadly;
the specific parameters are domain-dependent.

\subsection{Objection 2: Mechanism Mixtures}

\textbf{Objection:} Real behaviors may involve multiple mechanisms simultaneously.
The taxonomy is overly simplistic.

\textbf{Response:} Axiom HF-1 acknowledges this by specifying ``primary'' mechanism.
The value of the taxonomy is in providing testable predictions. Even if multiple
mechanisms operate, the temporal pattern reveals which dominates. Mixed models
can be estimated but require richer data.

\subsection{Objection 3: Selection Effects}

\textbf{Objection:} People who respond to first-period incentives may differ
systematically from those who don't, confounding the habit effect.

\textbf{Response:} This is why instrumental variables identification (using random
incentive assignment) is essential. Jones et al. addresses this through
rerandomization. The key is separating selection from true habit formation.

\subsection{Identification Boundary: First Contact vs. Quality of First Contact}

\begin{tcolorbox}[colback=red!5!white, colframe=red!75!black,
    title=Critical Limitation: What the Study Does NOT Identify]

\textbf{The study identifies:}
\begin{itemize}[nosep]
\item The causal effect of \textit{first participation vs. non-participation}
\item Formally: $\theta_1 = \mathbb{E}[S_t \mid \text{first contact induced by incentive}]$
\item This is an \textbf{extensive margin} effect
\end{itemize}

\textbf{The study does NOT identify:}
\begin{itemize}[nosep]
\item Differences in perceived quality of first experience
\item Effect of positive vs. negative first experience
\item Heterogeneity in experience outcomes
\item Formally: $\mathbb{E}[S_t \mid \text{good experience}] - \mathbb{E}[S_t \mid \text{bad experience}]$
\end{itemize}

\textbf{Why not?} Quality of first contact is:
\begin{itemize}[nosep]
\item \textbf{Endogenous} (not randomly assigned)
\item \textbf{Idiosyncratic} (varies by individual)
\item \textbf{Potentially correlated with types} (unobserved heterogeneity)
\item $\Rightarrow$ Not instrumentable with available data
\end{itemize}

\end{tcolorbox}

\textbf{Implicit model assumption (not tested):}
The paper implicitly assumes that first experience provides a sufficiently strong signal
about utility, \textit{regardless of idiosyncratic quality differences}. This assumption is:
\begin{itemize}[nosep]
\item Plausible
\item Consistent with experience-good models
\item But not empirically tested
\end{itemize}

\begin{tcolorbox}[colback=blue!5!white, colframe=blue!75!black,
    title=Correct Statements About Quality]

\textbf{Incorrect:} ``Quality doesn't matter.''

\textbf{Correct:} ``Regardless of (unobserved) quality, no additional learning effect emerges from subsequent experiences.''

\textbf{Precise formulation:}
\begin{quote}
\textit{``The study permits statements about the effect of first contact as such,
but not about the influence of the quality of that first experience.''}
\end{quote}

\textbf{Model-oriented:}
\begin{quote}
\textit{``The paper supports a discrete learning model at first contact,
but abstracts from heterogeneity in signal quality.''}
\end{quote}

\end{tcolorbox}


%%%%%%%%%%%%%%%%%%%%%%%%%%%%%%%%%%%%%%%%%%%%%%%%%%%%%%%%%%%%%%%%%%%%%%%%%%%%%%%
\section{Open Issues and Research Agenda}
\label{sec:hf-open}
%%%%%%%%%%%%%%%%%%%%%%%%%%%%%%%%%%%%%%%%%%%%%%%%%%%%%%%%%%%%%%%%%%%%%%%%%%%%%%%

\begin{table}[htbp]
\centering
\caption{Research Roadmap for Habit Formation Mechanisms}
\label{tab:hf-roadmap}
\renewcommand{\arraystretch}{1.3}
\begin{tabular}{@{}clcl@{}}
\toprule
\textbf{\#} & \textbf{Issue} & \textbf{Priority} & \textbf{Status} \\
\midrule
1 & Cross-domain replication of mechanism tests & HIGH & Open \\
2 & Interaction between mechanisms and heterogeneity & HIGH & Open \\
3 & Optimal incentive timing with mechanism uncertainty & MEDIUM & Open \\
4 & Long-run decay (beyond 3 years) & MEDIUM & Open \\
5 & Mechanism-specific experience quality interventions & HIGH & Open \\
\bottomrule
\end{tabular}
\end{table}

\subsection{Priority 1: Cross-Domain Replication}

The Jones et al. design should be replicated in other domains:
\begin{itemize}
\item Vaccination uptake (likely learning)
\item Exercise frequency (likely mixed)
\item Digital engagement (likely addiction)
\item Financial behaviors (likely inertia)
\end{itemize}

\subsection{Priority 5: Experience Quality Interventions}

If learning mechanism dominates, first-experience quality becomes a key lever.
Research needed on:
\begin{itemize}
\item What features of first experience affect permanent beliefs?
\item Can negative first experiences be ``overwritten''?
\item How does information provision interact with direct experience?
\end{itemize}


%%%%%%%%%%%%%%%%%%%%%%%%%%%%%%%%%%%%%%%%%%%%%%%%%%%%%%%%%%%%%%%%%%%%%%%%%%%%%%%
\section{References}
\label{sec:hf-references}
%%%%%%%%%%%%%%%%%%%%%%%%%%%%%%%%%%%%%%%%%%%%%%%%%%%%%%%%%%%%%%%%%%%%%%%%%%%%%%%

\begin{tcolorbox}[colback=blue!5!white, colframe=blue!75!black]
\textbf{Master Bibliography:} For complete reference details, see
\texttt{bcm2\_master\_references.bib} in the repository.

This section lists appendix-specific references and data sources.
\end{tcolorbox}

\vspace{0.5em}

\subsection{Key References for This Appendix}

The following works are particularly relevant for the content of this appendix:

\begin{itemize}[nosep]
    \item \textbf{Core Empirical:} Jones, Molitor \& Reif (2025)
    \item \textbf{Addiction Models:} Becker \& Murphy (1988)
    \item \textbf{Experience Goods:} Nelson (1970)
    \item \textbf{Related Habit Studies:} Hussam, Rabbani, Rigol \& Roth (2022); Charness \& Gneezy (2009)
\end{itemize}

\subsection{Appendix-Specific References}

\begin{thebibliography}{99}

\bibitem{jonesmolitorreif2025} Jones, D., Molitor, D., \& Reif, J. (2025).
How Do Habits Form? Experimental Evidence from Health Screenings.
\textit{Working Paper}.

\bibitem{beckermurphy1988} Becker, G. S., \& Murphy, K. M. (1988).
A Theory of Rational Addiction. \textit{Journal of Political Economy}, 96(4), 675-700.

\bibitem{nelson1970} Nelson, P. (1970).
Information and Consumer Behavior. \textit{Journal of Political Economy}, 78(2), 311-329.

\bibitem{hussam2022} Hussam, R., Rabbani, A., Rigol, N., \& Roth, B. (2022).
Habit Formation and Rational Addiction: A Field Experiment in Handwashing.
\textit{Working Paper}.

\bibitem{charnessgneezy2009} Charness, G., \& Gneezy, U. (2009).
Incentives to Exercise. \textit{Econometrica}, 77(3), 909-931.

\bibitem{madrianshea2001} Madrian, B. C., \& Shea, D. F. (2001).
The Power of Suggestion: Inertia in 401(k) Participation and Savings Behavior.
\textit{Quarterly Journal of Economics}, 116(4), 1149-1187.

\end{thebibliography}


%%%%%%%%%%%%%%%%%%%%%%%%%%%%%%%%%%%%%%%%%%%%%%%%%%%%%%%%%%%%%%%%%%%%%%%%%%%%%%%
% CROSS-REFERENCE FOOTER (Required)
%%%%%%%%%%%%%%%%%%%%%%%%%%%%%%%%%%%%%%%%%%%%%%%%%%%%%%%%%%%%%%%%%%%%%%%%%%%%%%%

\vspace{1em}
\hrule
\vspace{0.5em}

\noindent\textit{Cross-references:}
Appendix G (REF-GLOSSARY),
Appendix AW (CORE-STAGE),
Appendix AU (CORE-AWARE),
Appendix BBB (CORE-WHERE),
Appendix S (PREDICT-MASTER),
Appendix XIX (LIT-LEARNING).

\noindent\textit{Master Bibliography:} \texttt{bcm2\_master\_references.bib}


% =============================================================================
% END OF APPENDIX HF
% =============================================================================
